% Generated by lean-to-paper from paradoxical/BakerReduction.lean (2026-07-11),
% covering its project-local import closure: Defs.lean, RCircuitSlice.lean,
% ExcursionBounds.lean, RunPigeonhole.lean, PigeonholeFiniteness.lean,
% RhinPolyBound.lean, BakerReduction.lean.
% (C) 2026 Ralf Stephan, in collaboration with Claude Code. CC0 1.0 Universal.
\documentclass[11pt]{amsart}
\usepackage{amsmath,amssymb,amsthm}
\usepackage[hidelinks]{hyperref}
\usepackage{booktabs}

\theoremstyle{plain}
\newtheorem{theorem}{Theorem}[section]
\newtheorem{proposition}[theorem]{Proposition}
\newtheorem{lemma}[theorem]{Lemma}
\newtheorem{corollary}[theorem]{Corollary}
\newtheorem{conjecture}[theorem]{Conjecture}
\theoremstyle{definition}
\newtheorem{definition}[theorem]{Definition}
\theoremstyle{remark}
\newtheorem{remark}[theorem]{Remark}

\newcommand{\N}{\mathbb{N}}
\newcommand{\Z}{\mathbb{Z}}
\newcommand{\Par}{\mathcal{P}}        % the paradoxical set
\newcommand{\ParR}[1]{\Par_{#1}}      % the R-circuit slice

% ---------------------------------------------------------------------------
% Terminology table (lean-to-paper; the source of every technical term used).
% Rows shared with CanonicalShape.tex are kept identical for stability.
%
% Lean identifier / concept   | literature term            | source            | notation
% ----------------------------+----------------------------+-------------------+----------
% CC.T                        | Collatz function           | [Roz25, eq. (1)]  | T
% CC.T_iter j n               | iterate of T               | [Roz25, Sec. 1]   | T^j(n)
% pbits j n                   | parity vector              | [Roz25, Def. 2.1] | V_j(n)
% CC.num_odd_steps j n        | number of odd terms        | [Roz25, Def. 2.1] | q_j(n), q
% CC.C j n = 3^q/2^j          | coefficient                | [Roz25, eq. (2)]  | C_j(n), C
% q/j                         | ones-ratio                 | [Roz25, Def. 2.1] | q/j
% (iterate sequence)          | the sequence Omega_j(n)    | [Roz25, Def. 4.1] | \Omega_j(n)
% RT.IsParadoxical j n        | paradoxical (sequence)     | [Roz25, Def. 1.1] | --
% run-superscript words       | binary word 1^a 0^e ...    | [Roz25, Def. 2.1] | 1^{a}0^{e}
% (block 1^a 0^e)             | circuit                    | [Roz25, App. A]   | (a, e)
%                             |   (extended usage; see footnote in Sec. 2)
% RCircuitSlice R j n         | (descriptive: paradoxical  | --                | \mathcal{P}_R
%                             |  pairs, parity vector a    |
%                             |  word of R circuits)       |
% RT.delayCol / d_T           | delay                      | [Roz25, Def. 5.2] | d_T(n)
% RT.maxExcursion             | maximum excursion          | [Roz25, Def. 5.2] | M_T(n)
% excWin j n                  | greatest term of Omega_j(n)| [Roz25, Def. 5.2] | M_j(n)
%                             |  (windowed max excursion)  |  (windowed)       |
% minWin j n                  | smallest term of Omega_j(n)| [Roz25, Sec. 4:   | m_j(n)
%                             |                            |  "smallest term"] |
% RT.rExp j                   | r = j/floor(j log2/log3)   | [Roz25, eq. (12)] | r(j)
% RT.H j = 1/(2^r(j)-3)       | (descriptive: the bound of | [Roz25, eq. (12)] | H(j)
%                             |  Corollary 4.4)            |                   |
% (harmonic mean h)           | harmonic mean of the odd   | [Roz25, Cor. 4.4] | h
%                             |  terms of Omega_{j-1}(n)   |                   |
% Rhin.logForm                | the form Lambda =          | [Roz25, Prop. 6.3]| \Lambda
%                             |  u0 + u1 log2 + u2 log3    |  = [Rhi87]        |
% Rhin.logHeight              | max(|u1|,|u2|)             | [Roz25, Prop. 6.3]| (spelled out)
% d_T(n) <= alpha log n       | delay bound                | [Roz25, Conj. 6.2]| --
% M_T(n) <= n^beta            | maximum excursion bound    | [Roz25, eq. (16)] | --
% (runs of odd/even steps)    | run; ascent/descent phase  | [Roz25, App. A]   | --
% (Baker-type input)          | linear forms in logarithms | [Roz25, Sec. 6]   | --
% ---------------------------------------------------------------------------

\begin{document}

\title[Finiteness of paradoxical sequences with bounded circuits]{A conditional finiteness theorem for paradoxical Collatz sequences with a bounded number of circuits}
\author{Ralf Stephan}
\thanks{In collaboration with Claude Code.  All results stated below are
machine-verified in Lean~4; see Appendix~\ref{app:lean}}
\date{July 11, 2026}

\begin{abstract}
For the Collatz function $T$, Rozier and Terracol call a finite sequence
$\Omega_j(n)=(n,T(n),\dots,T^j(n))$ with $n>2$ \emph{paradoxical} when its
coefficient $C_j(n)=3^{q}/2^{j}$ is less than~$1$ and yet $T^j(n)\ge n$.  Fix
$R\ge 1$ and let $\ParR{R}$ be the set of paradoxical sequences whose parity
vector is a word of $R$ circuits $1^{a_1}0^{e_1}\cdots 1^{a_R}0^{e_R}$.  We
prove: if, for some $\beta\ge 0$, every member of $\ParR{R}$ satisfies the
excursion bound $M_j(n)\le m_j(n)^{\beta}$ --- where $M_j(n)$ and $m_j(n)$
denote the greatest and smallest terms of $\Omega_j(n)$ --- then $\ParR{R}$
is finite.  The two remaining links are unconditional: a run-length
pigeonhole argument shows $2^{j/(2R)}-1\le M_j(n)$ on the slice, and an
application of Rhin's effective irrationality measure for $\log 2/\log 3$ to
the denominator of Corollary~4.4 of Rozier--Terracol yields the polynomial
bound $m_j(n)\le\tfrac13 j^{14.3}$; an exponential--polynomial collision then
caps the length $j$.  The only transcendence input is Rhin's bound; all
results are machine-verified in Lean~4.
\end{abstract}

\maketitle

\section{Introduction}\label{sec:intro}

Among the necessary conditions that Rozier and Terracol \cite{Roz25} derive
for a Collatz sequence to be paradoxical, two have a quantitative character:
the ones-ratio of the parity vector is pinned close to $\log 2/\log 3$
\cite[Cor.~4.3]{Roz25}, and the harmonic mean of the odd terms is bounded
polynomially in the length \cite[Cor.~4.4]{Roz25}.  In \S6 of \cite{Roz25}
these are combined, through the effective irrationality measure of Rhin for
linear forms in $\log 2$ and $\log 3$ \cite{Rhi87} and conditionally on a
heuristic delay bound (\cite[Conj.~6.2]{Roz25}), into an argument that only
finitely many paradoxical sequences should exist.

This note isolates the part of that programme that can be carried out
rigorously for sequences with a \emph{bounded number of circuits}.  Writing
the parity vector of a paradoxical sequence as a word of $R$ circuits
$1^{a_1}0^{e_1}\cdots 1^{a_R}0^{e_R}$, we assemble three links:
\begin{enumerate}
\item[(1)] a \emph{run-length pigeonhole} (unconditional): some circuit
carries a run of one parity of length at least $j/(2R)$, and a long run of
either parity forces a large term, so the greatest term of the window
satisfies $2^{j/(2R)}-1\le M_j(n)$ (\S\S\ref{sec:runs}--\ref{sec:pigeonhole});
\item[(2)] a \emph{polynomial bound on the smallest term} (unconditional
modulo Rhin's theorem): $m_j(n)\le\tfrac13 j^{14.3}$
(\S\ref{sec:rhin});
\item[(3)] an \emph{excursion hypothesis} (the sole conjectural link): every
member of the slice satisfies $M_j(n)\le m_j(n)^{\beta}$ for a fixed
$\beta\ge 0$ --- a windowed form of the bound $M_T(n)\le n^{\beta}$ of
\cite[eq.~(16)]{Roz25}, itself a consequence of \cite[Conj.~6.2]{Roz25}.
\end{enumerate}
Chaining (1)--(3) traps an exponential in $j$ below a polynomial in $j$,
which bounds $j$; since a paradoxical sequence of length $j$ has first term
at most $2^j 3^j$, the slice is finite (Theorem~\ref{thm:main}).  A bound
uniform over \emph{all} $R$ cannot be expected from this route: it is
equivalent to the finiteness of the whole paradoxical set, hence implies the
Collatz conjecture (Remark~\ref{rem:uniform}).

\section{Notation and recalled results}\label{sec:background}

We follow the notation of \cite{Roz25}.  Let $T$ be the Collatz function
\[
T(n)=
\begin{cases}
\dfrac{3n+1}{2}, & n \text{ odd},\\[1ex]
\dfrac{n}{2}, & n \text{ even},
\end{cases}
\]
on the positive integers, $\Omega_j(n)=\bigl(T^k(n)\bigr)_{k=0}^{j}$ the
sequence of iterates of length $j+1$ \cite[Def.~4.1]{Roz25}, $q_j(n)$ (or
$q$) the number of odd terms among $n,T(n),\dots,T^{j-1}(n)$, the
\emph{coefficient} $C_j(n)=3^{q}/2^{j}$ \cite[eq.~(2)]{Roz25}, and
$V_j(n)\in\{0,1\}^j$ the \emph{parity vector} of $n$
\cite[Def.~2.1]{Roz25}, written as a binary word with repetitions
abbreviated by superscripts.

\begin{definition}[{\cite[Def.~1.1]{Roz25}}]\label{def:paradoxical}
For positive integers $j$ and $n$ with $n>2$, the sequence $\Omega_j(n)$ is
\emph{paradoxical} if
\begin{enumerate}
\item[(i)] $C_j(n)<1$, and
\item[(ii)] $T^j(n)\ge n$.
\end{enumerate}
We write $\Par=\{(j,n):\Omega_j(n)\text{ paradoxical},\ n>2\}$ for the
\emph{paradoxical set}.  Several intermediate statements below assume only
conditions (i)--(ii), with $n\ge 1$; we then say that $\Omega_j(n)$
satisfies (i)--(ii).
\end{definition}

Any $\Omega_j(n)$ satisfying (i)--(ii) with $n\ge1$ has $j\ge 2$: for $j=0$
one has $C_0(n)=1$, and for $j=1$ either $n$ is odd and $C_1(n)=3/2$, or $n$
is even and $T(n)=n/2<n$.

As in the companion note on canonical circuit decompositions, a
\emph{circuit} is a pair $(a,e)$ of nonnegative integers, spelling the word
$1^{a}0^{e}$; the terminology extends the classical usage, in which a
circuit is a (near-)cyclic sequence with a single ascent--descent phase,
i.e.\ with parity vector one block $1^{q}0^{j-q}$
\cite[App.~A]{Roz25}.\footnote{We are not aware of standard terminology for
the individual blocks of a general parity vector and adopt \emph{circuit}
for them.}  A finite list $S=\bigl((a_1,e_1),\dots,(a_R,e_R)\bigr)$ of
circuits spells the word
$w(S)=1^{a_1}0^{e_1}\cdots 1^{a_R}0^{e_R}$ of length
$|S|:=\sum_{i}(a_i+e_i)$.

\begin{definition}\label{def:slice}
For $R\ge 0$, the \emph{$R$-circuit slice} $\ParR{R}\subseteq\Par$ is the
set of pairs $(j,n)\in\Par$ such that $V_j(n)=w(S)$ for some list $S$ of $R$
circuits with $|S|=j$.  (Blocks may be empty, so the slices overlap; only
their union matters below.)
\end{definition}

Following \cite[Def.~5.2]{Roz25}, the \emph{delay} $d_T(n)$ is the least $j$
with $T^j(n)=1$ (or $+\infty$), and the \emph{maximum excursion} $M_T(n)$ is
the greatest term of the infinite sequence $\bigl(T^k(n)\bigr)_{k\ge 0}$ (or
$+\infty$).  We shall use their windowed counterparts
\[
M_j(n)=\max_{0\le k\le j}T^k(n),
\qquad
m_j(n)=\min_{0\le k\le j}T^k(n),
\]
the greatest and smallest terms of $\Omega_j(n)$; thus
$M_T(n)=\sup_j M_j(n)$.  Heuristic and numerical evidence suggests that the
delay satisfies $d_T(n)\le\alpha\log n$ for an absolute constant $\alpha$
\cite[Conj.~6.2]{Roz25}, from which one derives a polynomial excursion bound
$M_T(n)\le n^{\beta}$ \cite[eq.~(16)]{Roz25}.

The transcendence input is the effective irrationality measure of Rhin, as
stated in \cite[Prop.~6.3]{Roz25}.

\begin{proposition}[Rhin \cite{Rhi87}; see {\cite[Prop.~6.3]{Roz25}}]\label{prop:rhin}
If $u_0,u_1,u_2$ are rational integers such that
$\max(|u_1|,|u_2|)\ge 2$, then the form
$\Lambda=u_0+u_1\log 2+u_2\log 3$ verifies
\[
|\Lambda| \;\ge\; \max(|u_1|,|u_2|)^{-13.3}.
\]
\end{proposition}

Finally we recall two facts from the companion formal development.  The
first is the smallest-term form of \cite[Cor.~4.4]{Roz25}: the original
bounds the harmonic mean $h$ of the odd terms of $\Omega_{j-1}(n)$, and
\cite{Roz25} note in \S4 that one often uses the weakened form in which $h$
is replaced by the smallest term.  Write, as in \cite[eq.~(12)]{Roz25},
\[
H(j) \;=\; \frac{1}{2^{r(j)}-3},
\qquad
r(j) \;=\; \frac{j}{\bigl\lfloor j\log 2/\log 3\bigr\rfloor}.
\]

\begin{proposition}[smallest-term form of {\cite[Cor.~4.4]{Roz25}}]\label{prop:cor44}
Let $j\ge 2$, $n\ge 1$, and let $m\ge 1$ be a common lower bound for the odd
terms among $n,T(n),\dots,T^{j-1}(n)$.  If $\Omega_j(n)$ satisfies
(i)--(ii), then $m\le H(j)$.
\end{proposition}

\begin{lemma}\label{lem:startle}
If $\Omega_j(n)$ is paradoxical with $n>2$, then $n\le 2^{j}3^{j}$.
\end{lemma}

\section{Long runs force large terms}\label{sec:runs}

The arithmetic core of link (1) is that a run of consecutive steps of one
parity forces a large orbit term: $e$ successive halvings require
$2^{e}\mid T^{p}(n)$, while $a$ successive odd steps require
$2^{a}\mid T^{p}(n)+1$.

\begin{lemma}\label{lem:evenrun}
Let $n,e,p\ge 0$.  If $T^{p+i}(n)$ is even for all $i<e$, then
$2^{e}\mid T^{p}(n)$.
\end{lemma}

\begin{proof}[Proof sketch]
Induction on $e$: each even step halves the term, so $T^{p}(n)=2\,T^{p+1}(n)$
and the inductive divisibility at $p+1$ doubles.
\end{proof}

\begin{lemma}\label{lem:oddrun}
Let $n,a,p\ge 0$.  If $T^{p+i}(n)$ is odd for all $i<a$, then
$2^{a}\mid T^{p}(n)+1$.
\end{lemma}

\begin{proof}[Proof sketch]
Induction on $a$.  An odd step $x\mapsto(3x+1)/2$ sends $x+1$ to
$\tfrac{3}{2}(x+1)$: with $x=T^{p}(n)$ one checks
$T^{p+1}(n)+1=3\bigl(\lfloor x/2\rfloor+1\bigr)$, so the inductive
divisibility $2^{a-1}\mid T^{p+1}(n)+1$ passes to $2^{a-1}\mid
\lfloor x/2\rfloor+1$ (as $2^{a-1}$ is coprime to $3$) and doubles to
$2^{a}\mid x+1$.
\end{proof}

In particular a run of $e$ even steps forces $T^{p}(n)\ge 2^{e}$ (for
$n\ge1$), a run of $a$ odd steps forces $T^{p}(n)\ge 2^{a}-1$, and the
maximum excursion of Definition~5.2 of \cite{Roz25} satisfies
$M_T(n)\ge 2^{e}$, respectively $M_T(n)\ge 2^{a}-1$.

\section{The run-length pigeonhole}\label{sec:pigeonhole}

Throughout this section $S$ is a nonempty list of $R$ circuits with total
length $j=|S|$, and $n\ge 1$ satisfies $V_j(n)=w(S)$.  Paradoxicality is not
needed here.  The $i$-th circuit $(a_i,e_i)$ of $S$ contributes to $V_j(n)$
a run of $a_i$ odd steps followed by a run of $e_i$ even steps, located at
position $\sum_{i'<i}(a_{i'}+e_{i'})$ of the orbit.

\begin{proposition}\label{prop:shapeexc}
There are $L,p\ge 0$ with
\[
j\le 2RL,
\qquad
p\le j,
\qquad
2^{L}-1\le T^{p}(n).
\]
\end{proposition}

\begin{proof}[Proof sketch]
By the pigeonhole principle some circuit $(a_i,e_i)$ has
$a_i+e_i\ge j/R$; let $L=\max(a_i,e_i)\ge j/(2R)$.  If the even run is the
larger, Lemma~\ref{lem:evenrun} at its starting position $p$ gives
$2^{L}\mid T^{p}(n)$, so $T^{p}(n)\ge 2^{L}$; if the odd run is the larger,
Lemma~\ref{lem:oddrun} gives $2^{L}\mid T^{p}(n)+1$, so
$T^{p}(n)\ge 2^{L}-1$.  In both cases $p$ lies inside the window,
$p\le j$.
\end{proof}

\begin{proposition}\label{prop:pigeonhole}
With $S$, $j=|S|$, $R$, $n$ as above,
\[
2^{\,j/(2R)}-1 \;\le\; M_j(n).
\]
\end{proposition}

\begin{proof}[Proof sketch]
Take $L,p$ from Proposition~\ref{prop:shapeexc}.  Then
$2^{j/(2R)}\le 2^{L}$ since $j\le 2RL$, and
$2^{L}-1\le T^{p}(n)\le M_j(n)$ since $p\le j$.
\end{proof}

\section{The Rhin polynomial bound on the smallest term}\label{sec:rhin}

Link (2) bounds $H(j)$, and with it the smallest odd term of a sequence
satisfying (i)--(ii), polynomially in $j$.

\begin{proposition}\label{prop:Hbound}
For $j\ge 2$,
\[
H(j) \;\le\; \frac{j^{14.3}}{3}.
\]
\end{proposition}

\begin{proof}[Proof sketch]
Put $Q=\lfloor j\log 2/\log 3\rfloor$, so $1\le Q<j$ (using $j\ge 2$ and
$\log 2/\log 3<1$), and
$\Lambda_0=j\log 2-Q\log 3\ge 0$ by the floor.  Applying
Proposition~\ref{prop:rhin} with $(u_0,u_1,u_2)=(0,j,-Q)$ --- of height
$\max(j,Q)=j\ge 2$ --- gives $\Lambda_0\ge j^{-13.3}$.  On the other hand
$2^{r(j)}=3\,e^{\Lambda_0/Q}$, so by $e^{x}\ge 1+x$,
\[
2^{r(j)}-3 \;\ge\; \frac{3\Lambda_0}{Q},
\qquad\text{whence}\qquad
H(j)\;\le\;\frac{Q}{3\Lambda_0}\;\le\;\frac{j}{3}\cdot j^{13.3}
=\frac{j^{14.3}}{3}. \qedhere
\]
\end{proof}

\begin{proposition}\label{prop:mbound}
Let $j,n,m$ be integers with $n\ge 1$, $m\ge 1$, such that every odd term
among $n,T(n),\dots,T^{j-1}(n)$ is at least $m$.  If $\Omega_j(n)$ satisfies
(i)--(ii), then
\[
m \;\le\; \tfrac{1}{3}\, j^{14.3}.
\]
\end{proposition}

\begin{proof}[Proof sketch]
$j\ge 2$ as noted in \S\ref{sec:background}; chain
Proposition~\ref{prop:cor44} with Proposition~\ref{prop:Hbound}.
\end{proof}

\begin{remark}\label{rem:route}
The bound of Proposition~\ref{prop:Hbound} does not appear in \cite{Roz25}:
in \S6 they apply Rhin's estimate to the form
$\Lambda=j\log 2-q\log 3$ attached to the sequence itself.  Since
Proposition~\ref{prop:rhin} holds for arbitrary integer coefficients, it may
instead be fed the denominator of \cite[eq.~(12)]{Roz25}, whose exponent
$r(j)$ involves the floor $\lfloor j\log 2/\log 3\rfloor$ rather than the
odd-step count $q$ of any particular sequence; this makes the bound on
$H(j)$ --- and hence on the smallest odd term --- a function of the length
alone.  The exponent is $14.3=13.3+1$.
\end{remark}

\section{The collision engine and the main theorem}\label{sec:main}

The analytic engine is elementary: an exponential trapped below a
polynomial bounds the argument.

\begin{lemma}\label{lem:collision}
Let $a>0$ and $K,Q\ge 0$.  There is a $J\in\N$ such that every $j\in\N$ with
$e^{aj}-1\le K j^{Q}$ satisfies $j\le J$.
\end{lemma}

\begin{proof}[Proof sketch]
$t^{Q}e^{-at}\to 0$ as $t\to\infty$, so beyond some $M$ one has
$t^{Q}e^{-at}<1/(K+1)$; take $J=\lceil M\rceil$.  For $j>J$ this yields
$(K+1)j^{Q}<e^{aj}$, while the assumed inequality (together with
$j^{Q}\ge 1$) gives $e^{aj}\le K j^{Q}+1\le(K+1)j^{Q}$, a contradiction.
\end{proof}

We state the finiteness engine for an abstract class of sequences, with the
excursion and smallest-term functions as parameters; it will be instantiated
on the $R$-circuit slice.

\begin{theorem}\label{thm:engine}
Let $R\ge 1$ be an integer, let $\beta,P_w\ge 0$ and $C>0$ be real
constants, and let $P$ be a set of pairs $(j,n)$ such that every
$(j,n)\in P$ satisfies (i)--(ii) with $n>2$.  Let
$\mathrm{Exc},\,\mathrm{m}\colon P\to\mathbb{R}$ be functions such that, for
all $(j,n)\in P$,
\begin{enumerate}
\item[(1)] $2^{\,j/(2R)}-1\le\mathrm{Exc}(j,n)$,
\item[(2)] $1\le\mathrm{m}(j,n)\le C\,j^{P_w}$, and
\item[(3)] $\mathrm{Exc}(j,n)\le\mathrm{m}(j,n)^{\beta}$.
\end{enumerate}
Then $P$ is finite.
\end{theorem}

\begin{proof}[Proof sketch]
Write $2^{j/(2R)}=e^{aj}$ with $a=\log 2/(2R)>0$.  Chaining (1), (3), (2)
gives
\[
e^{aj}-1
\;\le\;\mathrm{Exc}(j,n)
\;\le\;\mathrm{m}(j,n)^{\beta}
\;\le\;\bigl(C j^{P_w}\bigr)^{\beta}
= C^{\beta} j^{P_w\beta},
\]
so Lemma~\ref{lem:collision} bounds $j\le J$ on $P$.  By
Lemma~\ref{lem:startle}, $n\le 2^{j}3^{j}\le 2^{J}3^{J}$, and $P$ lies in a
finite rectangle.
\end{proof}

Instantiating on the slice yields the main theorem.  Its sole hypothesis is
the excursion bound; the other two links are
Proposition~\ref{prop:pigeonhole} and Proposition~\ref{prop:mbound}.

\begin{theorem}\label{thm:main}
Let $R\ge 1$ be an integer and $\beta\ge 0$ a real constant.  Assume that
every $(j,n)\in\ParR{R}$ satisfies the excursion bound
\[
M_j(n) \;\le\; m_j(n)^{\beta}.
\]
Then $\ParR{R}$ is finite.
\end{theorem}

\begin{proof}[Proof sketch]
Apply Theorem~\ref{thm:engine} with $P=\ParR{R}$,
$\mathrm{Exc}=M_j(n)$, $\mathrm{m}=m_j(n)$, $C=\tfrac13$ and $P_w=14.3$.
Hypothesis (1) is Proposition~\ref{prop:pigeonhole} applied to the
witnessing circuit list of Definition~\ref{def:slice}.  For (2): all terms
of $\Omega_j(n)$ are at least $1$, so $1\le m_j(n)$; and since the smallest
term of the window is a common lower bound for its odd terms,
Proposition~\ref{prop:mbound} gives $m_j(n)\le\tfrac13 j^{14.3}$.
Hypothesis (3) is the assumption.
\end{proof}

\begin{remark}\label{rem:hexc}
The hypothesis of Theorem~\ref{thm:main} is a windowed form of the
polynomial excursion bound $M_T(n)\le n^{\beta}$ of \cite[eq.~(16)]{Roz25}
(a consequence of the delay bound \cite[Conj.~6.2]{Roz25}), with the
starting value replaced by the smallest term of the window.  Numerical
experiments on the computed portion of the slice are consistent with a
value $\beta\in[1.58,\,2.06]$.
\end{remark}

\begin{remark}\label{rem:uniform}
The per-$R$ statement is the honest endpoint of this route.  In the
companion note on canonical circuit decompositions it is shown that a
length bound valid on $\ParR{R}$ uniformly in $R$ is \emph{equivalent} to
the finiteness of the whole paradoxical set, which by
\cite[Cor.~3.3]{Roz25} implies the Collatz conjecture; and assuming the
excursion hypothesis uniformly in $R$ does not produce such a bound through
the present argument, since for $R\approx j/2$ the pigeonhole exponent
$j/(2R)$ degenerates to a constant.
\end{remark}

\appendix

\section{Formal statements}\label{app:lean}

{\sloppy
The results of this note are machine-verified in Lean~4.  The assembly
lives in \texttt{BakerReduction.lean}, in directory \texttt{paradoxical/}
of the author's formal development, together with the files it builds upon:
\texttt{Defs}, \texttt{RCircuitSlice}, \texttt{ExcursionBounds},
\texttt{RunPigeonhole}, \texttt{PigeonholeFiniteness},
\texttt{RhinPolyBound}.\par}  All proofs are
complete; the sole extra-logical axiom in the development is Rhin's theorem
(Proposition~\ref{prop:rhin}), recorded as a cited axiom on the authority of
\cite{Rhi87}.  Its phrasing was checked against the verbatim restatement
\cite[Prop.~6.3]{Roz25}; the original \cite{Rhi87} was not available to this
write-up.

Declarations live in the Lean namespace \texttt{Paradoxical} unless
prefixed otherwise.

\begin{center}
\footnotesize
\begin{tabular}{lll}
\toprule
Environment & Lean declaration & status \\
\midrule
Def.~\ref{def:paradoxical} & \texttt{RT.IsParadoxical} & definition (companion) \\
Def.~\ref{def:slice} & \texttt{RCircuitSlice} & definition \\
$M_j(n)$, $m_j(n)$ & \texttt{excWin}, \texttt{minWin} & definitions \\
$H(j)$, $r(j)$ & \texttt{RT.H}, \texttt{RT.rExp} & definitions (companion) \\
Prop.~\ref{prop:rhin} & \texttt{Rhin.logForm\_lower\_bound} & cited axiom [Rhi87] \\
Prop.~\ref{prop:cor44} & \texttt{RT.m\_le\_H} & proved (companion) \\
Lem.~\ref{lem:startle} & \texttt{RT.paradoxical\_start\_le} & proved (companion) \\
Lem.~\ref{lem:evenrun} & \texttt{two\_pow\_dvd\_of\_even\_run} & proved \\
Lem.~\ref{lem:oddrun} & \texttt{two\_pow\_dvd\_succ\_of\_odd\_run} & proved \\
Prop.~\ref{prop:shapeexc} & \texttt{shape\_excursion\_bound} & proved \\
Prop.~\ref{prop:pigeonhole} & \texttt{pigeonhole\_hyp} & proved \\
Prop.~\ref{prop:Hbound} & \texttt{H\_le\_rpow} & proved, modulo Rhin \\
Prop.~\ref{prop:mbound} & \texttt{least\_odd\_poly\_bound} & proved, modulo Rhin \\
Lem.~\ref{lem:collision} & \texttt{exp\_poly\_collision} & proved \\
Thm.~\ref{thm:engine} & \texttt{finite\_of\_pigeonhole} & proved \\
Thm.~\ref{thm:main} & \texttt{finite\_RCircuitSlice} & proved, modulo Rhin \\
\bottomrule
\end{tabular}
\end{center}

% omitted from the exposition (folded into prose or plumbing):
%   Defs.lean: remStep, ones, wordOfShape_append, wordOfShape_length,
%     wordOfShape_cons, wordOfShape_split, wlen_append, wlen_cons,
%     pbits_getElem?, odd_of_X_eq_one, even_of_X_eq_zero,
%     odd_of_word_bit_one, even_of_word_bit_zero
%   ExcursionBounds.lean: pow_le_maxExcursion_of_even_run,
%     pow_pred_le_maxExcursion_of_odd_run  (stated in prose, end of Sec. 3)
%   RunPigeonhole.lean: burst_odd_run, gap_even_run (position bookkeeping,
%     stated in the preamble of Sec. 4), exists_ge_avg (pigeonhole step)
%   BakerReduction.lean: one_le_minWin, minWin_le (window bounds, folded
%     into the proof sketch of Thm. 6.3)

\begin{thebibliography}{9}

\bibitem{Rhi87}
G.~Rhin,
\emph{Approximants de Pad\'e et mesures effectives d'irrationalit\'e},
S\'eminaire de Th\'eorie des Nombres, Paris 1985--86, 155--164,
Progr. Math. 71, Birkh\"auser Boston, Boston, MA, 1987.

\bibitem{Roz25}
O.~Rozier and C.~Terracol,
\emph{Paradoxical behavior in Collatz sequences},
arXiv:2502.00948 (2025).

\end{thebibliography}

\end{document}
